\section{Experimental Methodology}
\label{sec:methodology}

\subsection{Metrics, Corpus, and Environments}
\label{sec:corpus}

Security effectiveness is measured on two levels. Controlled penetration testing verifies each mechanism against a fixed suite of 63 test cases in 12 categories aligned with the OWASP LLM Top~10~\cite{owasp2025} and the NIST taxonomy of adversarial machine learning~\cite{nist2025aml}; every test specifies attack vector, targeted goal, expected layer, payload, and pass/fail criterion. Live evaluation submits a corpus of 42 attack queries in eight categories of increasing sophistication, naive extraction (8), instruction override (5), jailbreak and persona change (7), encoded payloads (4), typoglycemia (4), social engineering (5), indirect injection through fake system markers (4), and multi-turn escalation (5), together with 23 legitimate queries in three length groups (10 short, 10 medium, 3 long), to a deployed agent. It measures the \emph{block rate}, the share of attack queries that do not obtain a substantive answer, the \emph{false-positive rate} (FPR), the share of legitimate queries that are blocked, and, in the new experiments, the \emph{attack success rate} (ASR), the share of attack queries whose final response leaks protected content, so that blocking and leakage are measured separately. In the live evaluation of Section~\ref{sec:live}, an attack counts as blocked when the request is rejected, the response is empty, or the response is a refusal; the new experiments refine this criterion with the per-request classification of Section~\ref{sec:multimodel-protocol}. Performance overhead is measured as per-layer latency at the 50th, 95th, and 99th percentiles (P50, P95, P99) and as throughput under load; deployability as monthly infrastructure cost on three providers. Two environments are used: the two-account AWS deployment of Section~\ref{sec:implementation}, with Amazon Nova Lite as the underlying model, for the live evaluation, the performance measurements, and the cost analysis; and the multi-account Docker Compose stack of Section~\ref{sec:local-env}, which runs the same agent and gatekeeper code against open-weight models served by Ollama, for the multi-model evaluation and the ablation study. Both use the test prompts of the corpus rather than a production prompt.

\subsection{Multi-Model Evaluation Protocol}
\label{sec:multimodel-protocol}

The multi-model evaluation answers RQ1 by holding every element of the pipeline fixed and varying only the underlying LLM, so that differences in outcome can be attributed to the model. Nine models from six families are evaluated, covering both deployment modes of the architecture. Amazon Nova Lite (\texttt{amazon.nova-lite-v1:0}) is invoked through Amazon Bedrock and serves as the reference model of the live evaluation. The eight open-weight models are served by Ollama 0.33.3 in the local environment, with the quantization published under each identifier: Llama~3.2 1B (\texttt{llama3.2:1b}, Q8\_0), Llama~3 8B (\texttt{llama3:8b}, Q4\_0), Gemma~3 4B (\texttt{gemma3:4b}), Gemma~4 E4B (\texttt{gemma4:e4b}, 4B effective parameters), Phi-4 Mini (\texttt{phi4-mini}, 3.8B), Mistral~7B (\texttt{mistral:7b}), Qwen~3.5 4B (\texttt{qwen3.5:4b}), and Qwen~3.5 9B (\texttt{qwen3.5:9b}, 9.6B), the last six at Q4\_K\_M; seven of them ran on a single laptop GPU (RTX 4060, 8~GB) and Gemma~4 E4B on an RTX 3090. All models run the same agent and gatekeeper code, and serving conditions are constant within each model, which is what the comparison across configurations requires.

The following factors are fixed across models. The base system prompt is the medium-length prompt of the corpus (about 600 characters), whose closing sentence is a minimal confidentiality instruction that forms part of the base prompt in every configuration; it is hardened with the production defensive suffix and the canary tokens of Section~\ref{sec:hardening}. The four demonstration tools were registered for all models except Gemma~3 4B and Llama~3 8B, whose runs were executed without tool registration. Decoding is deterministic (temperature 0, nucleus sampling disabled, at most 1,024 output tokens on Ollama). All four layers are enabled, which corresponds to configuration C4 of Table~\ref{tab:ablation-configs}. Every query is submitted in a fresh conversation, except the multi-turn escalation category, whose five queries form one session and are submitted in their listed order. Each model is evaluated in three independent runs, reported as mean and standard deviation. Requests that end in a transport error or a timeout are counted as errors and excluded from the denominators; the complete campaign comprises 10,530 requests, 7 of which ended in errors.

The outcome of every request is classified by the mechanism that determined it: \emph{L2 block} (risk score at or above the threshold), \emph{model refusal} (the model declines before any filtering), \emph{L3 filter} (the response filter replaced the output), \emph{canary} (a canary token appeared in the output), \emph{pass} (a substantive response reached the caller), or \emph{error}. Because Layer~2 blocks, canary detections, and model refusals return the same refusal sentence, the label is recorded by the agent as a structured audit field rather than inferred from the response text; the raw model output produced before Layer~3 is recorded as well. Leakage is assessed offline with the same criterion that Layer~3 applies at runtime: a response leaks if it contains a canary token or reproduces any 50-character chunk of the base prompt, ignoring case. The oracle is applied both to the final response, which yields the ASR, and to the raw model output, which yields the \emph{pre-filter leakage} of the model. The block rate is the share of attack queries whose outcome is an L2 block, a model refusal, an L3 filter, or a canary detection; the FPR is the share of legitimate queries with any of these outcomes. Results are reported per model, with the attribution of outcomes to mechanisms (Table~\ref{tab:multimodel-overall}), and per attack category (Table~\ref{tab:multimodel-percategory}).

\subsection{Defense-Layer Ablation Protocol}
\label{sec:ablation-protocol}

\begin{table*}[t]
\caption{\label{tab:ablation-configs}Ablation configurations. Layer~1 (gatekeeper delivery) and Layer~4 (monitoring) remain active in every configuration, since they do not act on the input--output path of a request; they are verified separately (Table~\ref{tab:l1-l4-verification}). DP denotes the defensive suffix of Section~\ref{sec:hardening}.}
\centering
\footnotesize
\setlength{\tabcolsep}{3pt}
\begin{tabular}{@{}llccccccc@{}}
\toprule
\textbf{Config.} & \textbf{Description} & \textbf{DP} & \textbf{L2 scoring} & \textbf{L2 blocking} & \textbf{Sanitization} & \textbf{L3 filter} & \textbf{Canary} & \textbf{Watermark} \\
\midrule
C0            & No defenses (model-only floor)             & --         & --         & --         & --         & --         & --         & -- \\
C1            & DP only                                    & \checkmark & --         & --         & --         & --         & --         & -- \\
C2            & DP + Layer~2                               & \checkmark & \checkmark & \checkmark & \checkmark & --         & --         & -- \\
C3            & DP + Layer~3                               & \checkmark & --         & --         & --         & \checkmark & \checkmark & \checkmark \\
C4            & DP + Layer~2 + Layer~3 (deployed pipeline) & \checkmark & \checkmark & \checkmark & \checkmark & \checkmark & \checkmark & \checkmark \\
C4$^{\dagger}$ & As C4 with Layer~2 in detect-only mode    & \checkmark & \checkmark & --         & \checkmark & \checkmark & \checkmark & \checkmark \\
\bottomrule
\end{tabular}
\end{table*}


The ablation study answers RQ2 by measuring how the security outcome changes when individual defenses are removed. The layers are not interchangeable classifiers: Layer~1 governs authentication, request integrity, account binding, and replay protection, and Layer~4 governs monitoring, auditability, and coordinated response; neither acts on the content of a request or a response, so removing them does not change the outcome of the attack corpus. Both are kept active in every configuration and verified separately, with attack cases of their own, in Table~\ref{tab:l1-l4-verification}. The ablation therefore varies the \emph{content-path protections}, the mechanisms that act on the input--output path of a request: the defensive suffix, Layer~2, and Layer~3. Layer~1 and Layer~4 are validated functionally rather than as factors of the ablation.

Table~\ref{tab:ablation-configs} summarizes the configurations. C0 disables every content defense and measures the floor set by the model itself, since the base prompt retains only its closing confidentiality sentence; C1 adds the defensive suffix alone; C2 and C3 add Layer~2 and Layer~3, respectively, on top of the suffix; C4 is the deployed pipeline and coincides with the condition of the multi-model evaluation; and C4$^{\dagger}$ keeps the pipeline of C4 but runs Layer~2 in detect-only mode, in which the risk score is computed and recorded while blocking is disabled, so that every attack reaches the model and Layer~3 and the verdicts of both layers are observed on the same request. Input sanitization and canary detection follow the Layer~2 and Layer~3 settings, respectively. All configurations are evaluated for the nine models with the fixed factors, repetitions, classification, and metrics of Section~\ref{sec:multimodel-protocol}. The contribution of a mechanism is the change in block rate and ASR between configurations that differ only in that mechanism: the defensive suffix from C0 to C1, Layer~2 from C1 to C2 and from C3 to C4, and Layer~3 from C1 to C3 and from C2 to C4. Complementarity is measured on C4$^{\dagger}$ by labeling each attack query with the verdict of Layer~2 (would block or pass) and with its downstream outcome (intercepted by Layer~3 or by a canary, refused by the model, or passed), which yields Table~\ref{tab:complementarity}; queries that pass Layer~2 and are intercepted by Layer~3 are the direct evidence of complementarity, and queries that pass both are evaluated with the leakage oracle.

Layer~1 is verified with deterministic cases executed against the gatekeeper of the multi-account environment, which runs the production code: a valid signed request must return the prompt with a valid integrity hash, and each violated property must produce the response that follows from the validation order of the gatekeeper without disclosing prompt content. Layer~4 is verified by triggering each event class: the gatekeeper-side classes (repeated authentication failures, replay, rate limiting) in the same environment, whose alert topic is read through a subscribed queue, and the agent-side classes (Layer~2 blocks, Layer~3 filters, canary detections) and the customer-side control-plane rule on the two-account cloud deployment, where delivery is read from the archive log of the central event bus and from the invocation metrics of the forwarding rules, because the alarms of Layer~4 are cloud metric alarms that the local environment does not emulate.
